\documentclass[11pt]{article}
\usepackage[margin=1in]{geometry}
\usepackage{amsmath, amssymb}
\usepackage{booktabs}
\usepackage{enumitem}

\setlength{\parindent}{0pt}
\setlength{\parskip}{6pt}

\title{Assignment 2}
\author{Minh Nguyen}
\date{}

\begin{document}
\maketitle

%==============================================================
\section*{Part (a): $4 \times 3$ Grid World Reachability and Probabilities}

Starting at state $(1,1)$ with $P((1,1)) = 1.0$, executing the action sequence
$[\text{Up}, \text{Up}, \text{Right}, \text{Right}, \text{Right}]$ yields the following
step-by-step probability distribution.

\subsection*{Step-by-Step State Distributions}

\textbf{Step 1: Action = Up}
\begin{itemize}[nosep]
  \item $\text{Up} \rightarrow (1,2)$ with probability $0.8$
  \item $\text{Left} \rightarrow$ hits wall $\rightarrow (1,1)$ with probability $0.1$
  \item $\text{Right} \rightarrow (2,1)$ with probability $0.1$
\end{itemize}

\textbf{Step 2: Action = Up}
\begin{itemize}[nosep]
  \item From $(1,2)$: $\text{Up} \to (1,3)$ ($0.64$), Left/Right into wall/obstacle $\to (1,2)$ ($0.16$)
  \item From $(1,1)$: $\text{Up} \to (1,2)$ ($0.08$), $\text{Left} \to (1,1)$ ($0.01$), $\text{Right} \to (2,1)$ ($0.01$)
  \item From $(2,1)$: Up into obstacle $\to (2,1)$ ($0.08$), $\text{Left} \to (1,1)$ ($0.01$), $\text{Right} \to (3,1)$ ($0.01$)
\end{itemize}
State distribution: $(1,3): 0.64$, $(1,2): 0.24$, $(2,1): 0.09$, $(1,1): 0.02$, $(3,1): 0.01$

\textbf{Step 3: Action = Right}

State distribution: $(2,3): 0.512$, $(1,2): 0.258$, $(1,3): 0.088$, $(3,1): 0.073$,
$(2,1): 0.034$, $(1,1): 0.026$, $(4,1): 0.008$, $(3,2): 0.001$

\textbf{Step 4: Action = Right}

State distribution: $(3,3): 0.4097$, $(1,2): 0.2178$, $(2,3): 0.1728$, $(4,1): 0.0656$,
$(1,3): 0.0346$, $(3,1): 0.0346$, $(1,1): 0.0284$, $(2,1): 0.0276$, $(3,2): 0.0073$, $(4,2): 0.0016$

\textbf{Step 5: Action = Right}

\begin{center}
\begin{tabular}{cccl}
\toprule
Grid Cell & Exact Fraction & Decimal Probability & Percentage \\
\midrule
$(4,3)$ & $\frac{4099}{12500}$   & $0.32792$ & $32.792\%$ (Goal) \\[2pt]
$(1,2)$ & $\frac{9027}{50000}$   & $0.18054$ & $18.054\%$ \\[2pt]
$(3,3)$ & $\frac{8997}{50000}$   & $0.17994$ & $17.994\%$ \\[2pt]
$(4,1)$ & $\frac{543}{6250}$     & $0.08688$ & $8.688\%$ \\[2pt]
$(2,3)$ & $\frac{389}{6250}$     & $0.06224$ & $6.224\%$ \\[2pt]
$(3,2)$ & $\frac{4443}{100000}$  & $0.04443$ & $4.443\%$ \\[2pt]
$(2,1)$ & $\frac{353}{12500}$    & $0.02824$ & $2.824\%$ \\[2pt]
$(3,1)$ & $\frac{2627}{100000}$  & $0.02627$ & $2.627\%$ \\[2pt]
$(1,3)$ & $\frac{631}{25000}$    & $0.02524$ & $2.524\%$ \\[2pt]
$(1,1)$ & $\frac{1231}{50000}$   & $0.02462$ & $2.462\%$ \\[2pt]
$(4,2)$ & $\frac{171}{12500}$    & $0.01368$ & $1.368\%$ (Penalty) \\
\bottomrule
\end{tabular}
\end{center}

Every non-obstacle cell in the grid can be reached, with the intended goal state $(4,3)$
having the highest probability ($32.792\%$).

%==============================================================
\section*{Part (b): MDP Transformations \& Linear Property}

\subsection*{i. Transforming $R(s, a, s')$ into $R(s, a)$}

Define the expected immediate reward function $R'(s, a)$ as:
\[
R'(s, a) = \sum_{s'} P(s' \mid s, a)\, R(s, a, s')
\]

\textbf{Proof of Correspondence.} The Bellman expectation equation for any fixed policy
$\pi$ under $R(s, a, s')$ is:
\[
V^\pi(s) = \sum_{a} \pi(a \mid s) \sum_{s'} P(s' \mid s, a) \left[ R(s, a, s') + \gamma V^\pi(s') \right]
\]
Distributing the sum:
\[
V^\pi(s) = \sum_{a} \pi(a \mid s) \left[ \underbrace{\sum_{s'} P(s' \mid s, a) R(s, a, s')}_{R'(s, a)} + \gamma \sum_{s'} P(s' \mid s, a) V^\pi(s') \right]
\]
\[
V^\pi(s) = \sum_{a} \pi(a \mid s) \left[ R'(s, a) + \gamma \sum_{s'} P(s' \mid s, a) V^\pi(s') \right]
\]
Because $V^\pi(s)$ is identical under both formulations for every policy $\pi$, the expected
return for all policies is preserved. Thus, $\pi^* = \arg\max_\pi V^\pi(s)$ is identical in
both MDPs.

\subsection*{ii. Transforming $R(s, a)$ into $R(s)$}

To remove action dependence from the reward function, augment the state space by
introducing intermediate state-action states:
\begin{itemize}
  \item Define a new state space $S' = S \cup (S \times A)$.
  \item When the agent in state $s \in S$ chooses action $a$, it transitions
        deterministically to the intermediate state $s_a = (s, a)$ with reward $R(s) = 0$.
  \item From state $s_a = (s, a)$, the environment automatically transitions to state
        $s' \in S$ according to $P(s' \mid s, a)$ with state reward $R(s_a) = R(s, a)$.
\end{itemize}
Under this construction, every policy in the original MDP maps to an equivalent policy in
the transformed MDP. The cumulative discounted rewards match up to a constant discount
factor scaling, preserving the exact ordering and optimal policies $\pi^*$.

\subsection*{iii. Proof that $V^\pi(s)$ Varies Linearly with $R(s)$}

Let $R$ be the column vector of state rewards $R(s)$, $V^\pi$ be the value vector for
policy $\pi$, and $P^\pi$ be the transition probability matrix under $\pi$. The Bellman
equation in vector form is:
\[
V^\pi = R + \gamma P^\pi V^\pi
\quad\Longrightarrow\quad
(I - \gamma P^\pi) V^\pi = R
\]
Since $\gamma < 1$ and $P^\pi$ is a row-stochastic matrix, $(I - \gamma P^\pi)$ is
invertible:
\[
V^\pi = (I - \gamma P^\pi)^{-1} R
\]
Let $M = (I - \gamma P^\pi)^{-1}$. For any individual state $s$:
\[
V^\pi(s) = \sum_{s'} M_{s, s'} R(s')
\]
This is a linear combination of the components of $R$. Thus, $V^\pi(s)$ varies linearly
with $R(s)$.

%==============================================================
\section*{Part (c): 3-State Undiscounted MDP}

\subsection*{i. Qualitative Analysis of Optimal Policy}

\textbf{State 1} ($R(1) = -1$): Action $b$ provides a direct path to the terminal state 3
with probability $0.1$. Action $a$ moves the agent to state 2, which incurs a harsher
penalty per step ($R(2) = -2$). Thus, state 1 should select action $b$ to exit as quickly
as possible.

\textbf{State 2} ($R(2) = -2$): Taking action $b$ gives an expected value of
$V(2) = -2 + 0.9 V(2) \implies V(2) = -20$. Taking action $a$ moves the agent with
probability $0.8$ to state 1, where cost per step is lower ($-1$) and escape is faster.
Thus, state 2 should select action $a$ to transfer to state 1.

\subsection*{ii. Policy Iteration starting with $\pi_0 = (b, b)$}

\textbf{Iteration 1}

\emph{Policy Evaluation of $\pi_0 = (b, b)$:}
\begin{align*}
V^{\pi_0}(1) &= -1 + 0.9 V^{\pi_0}(1) + 0.1(0) \implies 0.1 V^{\pi_0}(1) = -1 \implies \mathbf{V^{\pi_0}(1) = -10} \\
V^{\pi_0}(2) &= -2 + 0.9 V^{\pi_0}(2) + 0.1(0) \implies 0.1 V^{\pi_0}(2) = -2 \implies \mathbf{V^{\pi_0}(2) = -20}
\end{align*}

\emph{Policy Improvement:}
\[
Q(s, a) = R(s) + \sum_{s'} P(s' \mid s, a)\, V^{\pi_0}(s')
\]
State 1:
\begin{align*}
Q(1, a) &= -1 + 0.2(-10) + 0.8(-20) = -19 \\
Q(1, b) &= -1 + 0.9(-10) + 0.1(0) = -10 \\
\pi_1(1) &= \arg\max \{-19, -10\} = b
\end{align*}
State 2:
\begin{align*}
Q(2, a) &= -2 + 0.8(-10) + 0.2(-20) = -14 \\
Q(2, b) &= -2 + 0.9(-20) + 0.1(0) = -20 \\
\pi_1(2) &= \arg\max \{-14, -20\} = a
\end{align*}
New policy: $\pi_1 = (b, a)$.

\textbf{Iteration 2}

\emph{Policy Evaluation of $\pi_1 = (b, a)$:}
\begin{align*}
V^{\pi_1}(1) &= -1 + 0.9 V^{\pi_1}(1) \implies \mathbf{V^{\pi_1}(1) = -10} \\
V^{\pi_1}(2) &= -2 + 0.8 V^{\pi_1}(1) + 0.2 V^{\pi_1}(2) \implies 0.8 V^{\pi_1}(2) = -2 + 0.8(-10) = -10 \implies \mathbf{V^{\pi_1}(2) = -12.5}
\end{align*}

\emph{Policy Improvement:}

State 1:
\begin{align*}
Q(1, a) &= -1 + 0.2(-10) + 0.8(-12.5) = -13 \\
Q(1, b) &= -1 + 0.9(-10) = -10 \\
\pi_2(1) &= \arg\max \{-13, -10\} = b
\end{align*}
State 2:
\begin{align*}
Q(2, a) &= -2 + 0.8(-10) + 0.2(-12.5) = -12.5 \\
Q(2, b) &= -2 + 0.9(-12.5) = -13.25 \\
\pi_2(2) &= \arg\max \{-12.5, -13.25\} = a
\end{align*}
Since $\pi_2 = \pi_1 = (b, a)$, Policy Iteration converges.

\begin{center}
\fbox{\parbox{0.7\textwidth}{\centering
Optimal Policy: $\pi^*(1) = b, \; \pi^*(2) = a$ \\[2pt]
Optimal Values: $V^*(1) = -10, \; V^*(2) = -12.5$}}
\end{center}

\subsection*{iii. Initial Policy $\pi_0 = (a, a)$, Discounting, and Dependency on $\gamma$}

\textbf{Initial Policy $\pi_0 = (a, a)$.} In the undiscounted case ($\gamma = 1$), the policy
$(a, a)$ cycles indefinitely between state 1 and state 2 without any probability of
reaching state 3. The matrix $(I - P)$ is singular, making $V(1)$ and $V(2)$ equal to
$-\infty$ (unbounded cost). Standard policy evaluation fails.

\textbf{Does Discounting Help?} Yes. For $\gamma < 1$, geometric discounting guarantees
convergence even for infinite trajectories:
\begin{align*}
V(1) &= -1 + \gamma\big(0.2 V(1) + 0.8 V(2)\big) \\
V(2) &= -2 + \gamma\big(0.8 V(1) + 0.2 V(2)\big)
\end{align*}
This yields finite values for $V(1)$ and $V(2)$, enabling the policy improvement step to
properly calculate $Q(s, b)$ and update the policy to $(b, a)$.

\textbf{Does the Optimal Policy Depend on $\gamma$?} No. For any $\gamma \in (0, 1]$:
\begin{align*}
V^*(1) &= \frac{-1}{1 - 0.9\gamma} \\
V^*(2) &= \frac{-2 + 0.6\gamma}{1 - 0.2\gamma - 0.72\gamma^2 + 0.18\gamma^2}
\end{align*}
Evaluating
\[
Q(1, b) - Q(1, a) = \frac{0.8\gamma(1 - 0.8\gamma)}{(1 - 0.2\gamma)(1 - 0.9\gamma)} > 0
\]
Evaluating
\[
Q(2, a) - Q(2, b) = \frac{0.6\gamma}{1 - 0.2\gamma} > 0
\]
The optimal policy remains $\pi^* = (b, a)$ for all valid discount factors $\gamma \in (0, 1]$.

%==============================================================
\section*{Part (d): Parametric Value Function Approximation \& TD Updates}

Given the value function approximator:
\[
\hat{V}_\theta(x, y) = \theta_0 + \theta_1 x + \theta_2 y + \theta_3 \sqrt{(x - x_g)^2 + (y - y_g)^2}
\]
The Temporal Difference (TD) error $\delta$ is defined as:
\[
\delta = R(x, y) + \gamma \hat{V}_\theta(x', y') - \hat{V}_\theta(x, y)
\]
Using the Widrow--Hoff gradient descent update rule:
\[
\theta_i \leftarrow \theta_i + \alpha\, \delta\, \frac{\partial \hat{V}_\theta(x, y)}{\partial \theta_i}
\]

\subsection*{Partial Derivatives and TD Update Equations}

For $\theta_0$:
\[
\frac{\partial \hat{V}_\theta(x, y)}{\partial \theta_0} = 1 \implies \boldsymbol{\theta_0 \leftarrow \theta_0 + \alpha \delta}
\]
For $\theta_1$:
\[
\frac{\partial \hat{V}_\theta(x, y)}{\partial \theta_1} = x \implies \boldsymbol{\theta_1 \leftarrow \theta_1 + \alpha \delta x}
\]
For $\theta_2$:
\[
\frac{\partial \hat{V}_\theta(x, y)}{\partial \theta_2} = y \implies \boldsymbol{\theta_2 \leftarrow \theta_2 + \alpha \delta y}
\]
For $\theta_3$:
\[
\frac{\partial \hat{V}_\theta(x, y)}{\partial \theta_3} = \sqrt{(x - x_g)^2 + (y - y_g)^2} \implies \boldsymbol{\theta_3 \leftarrow \theta_3 + \alpha \delta \sqrt{(x - x_g)^2 + (y - y_g)^2}}
\]

\end{document}